\scanpage{56}
Si
\[
\alpha:\Delta_{r'}\longrightarrow\Delta_r,\qquad \uuline{$\alpha(\unsure{0|\omega})=0$},
\]
\[
(29)\qquad\left\{\begin{aligned}
&\text{alors}\qquad \alpha^*:X_{\Delta_r}\mathbin{\unsure{-|\longrightarrow}}X_{\Delta_{r'}}\qquad(\text{étale})\\
&\text{satisfait}\\
&\qquad\alpha^*(X_{\Delta_{r'}}^*)=X_{\Delta_r}^*,\quad\text{ie. }\alpha^*(\dot X_{\Delta_{r'}})=\dot X_{\Delta_r}.
\end{aligned}\right.
\]

\deleted{\unsure{Dans le cas}}

\uline{NB} ``Moralement'', $X_{\Delta_r}^*$ représente une sorte de partie ``non singulière'' de $X_{\Delta_r}$, $\dot X_{\Delta_r}$ la partie singulière ou ``bord''. On fera attention que les $X_{\Delta_r}^*$ ne \uline{forment} \uline{pas} un espace semi-simplicial.

Pourtant, notre propos est, dans une certaine mesure, de reconstituer $X$ à partir des ``strates ouvertes'' $X_i^*$ dans $X$, dans l'esprit du n° 1 (où il y a deux strates $X_0=Y$ et $X_1=X$, \deleted{$\unsure{Y\subset X}$ \unsure{fermé}} \inserted{($Y\subset X$ inclusion} fermée), et on a donc $Y=Y^*$, $X^*=X\setminus Y$ comme dans les notations du loc. cit, où il y avait lieu d'introduire $\mathcal{V}_{\unsure{Y},X}$.
\margin{55}
